\documentclass[12pt,a4paper]{amsart}

\usepackage{amsmath,amssymb,amsthm}
\usepackage{mathtools}
\usepackage{enumitem}
\usepackage{hyperref}
\usepackage{booktabs}
\usepackage{tikz-cd}
\usetikzlibrary{positioning,arrows.meta}

%%%─ Environments
\newtheorem{theorem}{Theorem}[section]
\newtheorem{lemma}[theorem]{Lemma}
\newtheorem{proposition}[theorem]{Proposition}
\newtheorem{corollary}[theorem]{Corollary}
\newtheorem{conjecture}[theorem]{Conjecture}
\theoremstyle{definition}
\newtheorem{definition}[theorem]{Definition}
\newtheorem{hypothesis}[theorem]{Hypothesis}
\newtheorem{remark}[theorem]{Remark}

%%%─ Macros
\newcommand{\cA}{\mathcal{A}}
\newcommand{\cAg}{\mathcal{A}_\gamma}
\newcommand{\R}{\mathbb{R}}
\newcommand{\N}{\mathbb{N}}
\newcommand{\cK}{K_{B_c}}
\newcommand{\cKA}{K_{\cA}}
\newcommand{\normh}[2]{\|#1\|_{#2}}
\newcommand{\SK}{S_K}
\newcommand{\AK}{A_K}
\newcommand{\AKlin}{A_K^{\mathrm{lin}}}
\newcommand{\AKquad}{A_K^{\mathrm{quad}}}
\newcommand{\BK}{B_K}
\newcommand{\Kopt}{K^*}
\newcommand{\Clin}{C_{\mathrm{lin}}(M,K)}
\newcommand{\Cquad}{C_{\mathrm{quad}}(M,K)}
\newcommand{\Ctwo}{C_2(\alpha)}
\newcommand{\Cinfty}{C_\infty}
%%% Central notation (must appear identically throughout):
% \Phi_1, \Phi_3, \Psi       -- truncation terms and their sum
% \mathcal{M}(\mathrm{S1,S2,S3})  -- method class
% \mathfrak{R}               -- realization map
% c\text{-uniform stability} -- stability property of \mathfrak{R}

\subjclass[2020]{34B24, 47B34, 34E20, 47G10, 34L20}
\keywords{Sturm--Liouville operators, turning-point asymptotics,
  spectral growth exponent, two-scale truncation, logarithmic rate,
  spectral-geometric invariant, universality, layer-survival,
  spectrally regularized WKB class, variational truncation,
  Langer error, discrete floor projection, iterated asymptotics,
  diagonal substitution, stability of optimal truncation,
  substitution stability principle, structural obstruction,
  rigidity of scaling regimes, decomposition stability,
  balance fixed-point, method-class invariant,
  H1 exponent rigidity, envelope conjecture,
  scale admissibility, EM structural property,
  asymptotic preorder, realization map, $c$-uniform stability}

\begin{document}

\title{Universality and Structural Rigidity of
  Turning-Point Truncation Geometry}

\author{Sebastian Schmalnauer}
\address{Vienna, Austria}
\email{sebastian.schmalnauer@gmx.at}
\date{May 2026}

\begin{abstract}
For operators $\cAg$ in the spectrally regularized WKB class
$\mathfrak{T}_\gamma$ (Definition~\ref{def:opclass}),
the optimal two-scale truncation satisfies
$K^*(c) = \frac{2\beta}{\alpha}\log c + \frac{2\gamma}{\alpha}\log\log c + O(1)$,
yielding the diagonal rate
$O(c^{-\beta}(\log c)^{1+\gamma})$
for every operator in $\mathfrak{T}_\gamma$ satisfying H1--H2
(Corollary~\ref{cor:universal_lograte}).

The three structural components play distinct, mutually irreducible roles:
(S1) generates the exponent pair $(\gamma,\beta)$;
(S2) --- the Euler--Maclaurin structural property ---
ensures $c$-uniform stability of the realization map
$\mathfrak{R}\colon \{\text{analytic bounds}\} \to \mathcal{S}/\!\sim_{\mathcal{S}}$
(Definition~\ref{def:realization_map}),
so that leading-order bounds project to fixed equivalence classes
under $c$-rescaling;
(S3) licenses the diagonal substitution $K\mapsto K^*(c)$.

The central result is structural rigidity
(Theorem~\ref{thm:rate_obstruction}):
within $\mathcal{M}(\mathrm{S1,S2,S3})$,
no asymptotically better rate is compatible with the two-scale geometry.
Rigidity rests on decomposition stability (Lemma~\ref{lem:decomp_stability}),
balance fixed-point uniqueness (Lemma~\ref{lem:balance_fixpoint}),
H1 exponent rigidity (Lemma~\ref{lem:H1_rigidity}),
and the S2 knockdown (Lemma~\ref{lem:S2_knockdown}).

H1 envelope minimality (O1, Conjecture~\ref{conj:H1_envelope})
is confirmed for the Airy case; the general case remains open.
All results are conditional on H1--H2.
\end{abstract}

\maketitle
\tableofcontents
\bigskip\hrule\bigskip

%% ============================================================
\section{Introduction and Logical Status}
\label{sec:intro}
%% ============================================================

%%% Status vocabulary (unified throughout):
%%%   theorem     -- unconditional within stated hypotheses
%%%   conditional -- depends on H1--H2 or \mathfrak{T}_\gamma
%%%   open        -- not established in this paper

\begin{center}
\renewcommand{\arraystretch}{1.3}
\begin{tabular}{l l l}
\toprule
\textbf{Paper} & \textbf{Content} & \textbf{Status} \\
\midrule
XVIII & BR3: qualitative Airy universality & theorem \\
XIX   & Two-scale rate; lin/quad split     & theorem (H1--H2 verified) \\
XX    & Universal rate; structural obstruction;
        H1 rigidity; S2 knockdown; Layer-Survival
      & conditional ($\mathfrak{T}_\gamma$, H1--H2) \\
\bottomrule
\end{tabular}
\end{center}

\begin{remark}[Roles of S1, S2, S3]
\label{rem:S123_roles}
The three structural components of Theorem~\ref{thm:rate_obstruction}
play logically distinct and mutually irreducible roles:
\begin{itemize}[nosep]
  \item \textbf{(S1) = H1 (Langer scaling):}
    generates the exponent pair $(\gamma,\beta)$ and produces
    the algebraic growth $\Phi_1 = C_{\rm lin}K^{1+\gamma}c^{-\beta}$.
    Without S1, no growth term exists and the balance equation is vacuous.
  \item \textbf{(S2) = EM structural property:}
    ensures $c$-uniform stability of the realization map $\mathfrak{R}$
    (Definition~\ref{def:realization_map}):
    each analytic bound projects to a fixed equivalence class in
    $\mathcal{S}/\!\sim_{\mathcal{S}}$ under $c$-rescaling.
    The preorder $\succ$ on $\mathcal{S}$ is defined purely on
    index pairs and is unconditionally well-defined;
    S2 protects the stability of the map from bounds to those index pairs,
    not the preorder itself.
    Without S2, $\mathfrak{R}$ is not $c$-uniformly stable,
    the comparison $\Phi_1 \succ \Phi_2$ as applied to actual bounds
    ceases to be $c$-uniform,
    and the balance fixed-point argument of
    Lemma~\ref{lem:balance_fixpoint}(iii) collapses
    (Lemma~\ref{lem:S2_knockdown}).
    EM is a canonical realization of this stability property
    in the WKB power-sum class; other summation methods satisfying
    leading-power extraction, uniform remainder control, and scaling
    invariance of coefficients serve equally.
  \item \textbf{(S3) = Substitution Stability:}
    licenses the passage from the fixed-parameter regime
    to the diagonal regime ($K = K^*(c)$).
    Without S3, the optimality condition~\eqref{eq:FOC}
    cannot be substituted into the truncation bound.
\end{itemize}
The three roles are:\quad
exponent generation (S1),\quad
$c$-uniform stability of $\mathfrak{R}$ (S2),\quad
diagonal projection (S3).
\end{remark}

\begin{remark}[Scope of XX and meaning of universality]
\label{rem:logstatus}
All results are conditional on H1--H2 within $\mathfrak{T}_\gamma$.
Universality has two layers:
\begin{itemize}[nosep]
  \item \textbf{Rate universality (conditional):}
    for each fixed $\gamma>0$, the rate $O(c^{-\beta}(\log c)^{1+\gamma})$
    holds uniformly over all operators in $\mathfrak{T}_\gamma$
    satisfying H1--H2.
  \item \textbf{Structural rigidity (conditional):}
    the exponent $1+\gamma$ and the scale $\log c$ are fixed by
    the interplay of S1--S3 within $\mathcal{M}(\mathrm{S1,S2,S3})$;
    no improvement is possible within this method class.
\end{itemize}
\end{remark}

\begin{remark}[$c$-admissible constants]
\label{rem:cadmissible}
A constant is \emph{$c$-admissible} if it is independent of $c$.
Every constant falls under one of four declared sources:
\begin{itemize}[nosep]
  \item \textbf{Axiomatic:} $C_0$, $C_\infty$ by H1--H2.
  \item \textbf{Analytic:} $C_\sigma = \frac{|B_2|}{12}\sigma(\sigma-1)$.
  \item \textbf{Derived:} products/sums of $c$-admissible constants
    (Lemma~\ref{lem:cadm_closure}).
  \item \textbf{Operator-intrinsic:} $C_2(\alpha)$.
\end{itemize}
\end{remark}

\begin{lemma}[$c$-admissible closure]
\label{lem:cadm_closure}
The class of $c$-admissible constants is closed under
finite products and finite sums.
\end{lemma}
\begin{proof}
$A,B$ independent of $c$ $\Rightarrow$ $AB$ and $A+B$ independent of $c$.
\end{proof}

\begin{remark}[Two sources of $c$-dependence, strictly separated]
\label{rem:dag_csplit}
$c$ appears in two logically distinct roles:
\begin{itemize}[nosep]
  \item \textbf{Operator estimates:} $c^{-\beta}$ as scaling factor;
    all multiplying constants $c$-admissible.
  \item \textbf{Optimization domain:} $c$ in window
    $[c_1\log c, c_2\log c]$ and $K^*(c)$ via FOC.
\end{itemize}
The formal bridge is Proposition~\ref{prop:substitution_stability}.
\end{remark}

\begin{remark}[Two-regime quantifier structure]
\label{rem:diagonal_regime}
\begin{itemize}[nosep]
  \item \textbf{Fixed-parameter:}
    $\forall K\geq K_0,\, c\geq c_0$ independently.
  \item \textbf{Diagonal:}
    $\forall c\geq c_0,\;\exists!\,K^*(c)\in[c_1\log c,c_2\log c]$.
\end{itemize}
Transition licensed by Proposition~\ref{prop:substitution_stability}.
\end{remark}

%% ============================================================
\section{Substitution Stability Principle}
\label{sec:substprin}
%% ============================================================

\begin{proposition}[Substitution Stability Principle]
\label{prop:substitution_stability}
\emph{Intent.}
Substituting a deterministic map $K(c)$ into a
pointwise-valid inequality is admissible as a pointwise composition.

\emph{Setup.}
Let $|E(K,c)| \leq \sum_{n=1}^N C_n\,K^{a_n}c^{-b_n}$
hold for all $K\geq K_0$, $c\geq c_0$, $c$-admissible $C_n$.
Let $K(c)\in[c_1\log c,\,c_2\log c]$ for all large $c$.

\begin{itemize}[nosep]
  \item[(i)] \textbf{Pointwise substitution.}
    $|E(K(c),c)|\leq\sum_n C_n(K(c))^{a_n}c^{-b_n}$.
  \item[(ii)] \textbf{Window control via monotonicity.}
    $\sum_n C_n (K(c))^{a_n} c^{-b_n}
      \leq \sum_n C_n (c_2\log c)^{a_n} c^{-b_n}$.
  \item[(iii)] \textbf{Asymptotic order.}
    $(K(c))^{a_n}=O(\log^{a_n}c)$;
    precise correction by~\eqref{eq:Kstar_power}.
  \item[(iv)] \textbf{Dominance.}
    By Propositions~\ref{prop:asymp_order} and~\ref{prop:scale_algebra}.
\end{itemize}
\end{proposition}

\begin{proof}
(i) pointwise; (ii) monotonicity + window upper bound;
(iii) sandwich~\eqref{eq:sandwich}; (iv) Props.~\ref{prop:asymp_order}--\ref{prop:scale_algebra}.
\end{proof}

%% ============================================================
\section{Asymptotic Scale, Preorder, and Realization Map}
\label{sec:asymp}
%% ============================================================

\begin{definition}[Asymptotic scale and preorder]
\label{def:scale}
For $s>0$, $t\in\R$, let $\phi_{s,t}(c) := c^{-s}(\log c)^t$.
Define the \emph{asymptotic scale}
$\mathcal{S} := \{\phi_{s,t} : s>0,\, t\in\R\}$.

Two functions $f,g:(c_0,\infty)\to(0,\infty)$ are
\emph{asymptotically equivalent in $\mathcal{S}$},
written $f\sim_{\mathcal{S}} g$, if
\[
  0 < \liminf_{c\to\infty} \frac{f(c)}{g(c)}
  \leq \limsup_{c\to\infty} \frac{f(c)}{g(c)} < \infty,
\]
where both ratio bounds are finite positive real numbers,
\emph{independent of all external parameters}
(no dependence on $c$, $K$, $\lambda$, or any auxiliary index).

The \emph{asymptotic preorder} $\succ$ on $\mathcal{S}$ is
\[
  \phi_{s_1,t_1} \succ \phi_{s_2,t_2}
  \;\iff\;
  s_1 < s_2 \quad \text{or} \quad (s_1 = s_2 \text{ and } t_1 > t_2).
\]
This is the lexicographic order on $(s,t)\in\R^2$, defined
purely on index pairs, independently of any function representation.
The preorder $\succ$ is unconditionally well-defined.
\end{definition}

\begin{definition}[Realization map $\mathfrak{R}$]
\label{def:realization_map}
Let $\mathcal{B}$ denote the class of analytic bounds of the form
$b(K,c) = A(K,c)\,\phi_{s,t}(c)$, where $A(K,c)$ is a
non-negative prefactor and $\phi_{s,t}\in\mathcal{S}$.

The \emph{realization map} is
\[
  \mathfrak{R}\colon \mathcal{B} \to \mathcal{S}/\!\sim_{\mathcal{S}},
  \qquad
  \mathfrak{R}(b) := [\phi_{s,t}]
  \quad\text{if}\quad b(K^*(c),c) \sim_{\mathcal{S}} \phi_{s,t}.
\]
The map $\mathfrak{R}$ is \emph{$c$-uniformly stable} if,
for every $\lambda>0$ and every $b\in\mathcal{B}$,
\[
  \mathfrak{R}(b(\,\cdot\,,\lambda c)) = \mathfrak{R}(b(\,\cdot\,,c)).
\]
That is: rescaling $c\mapsto\lambda c$ does not change the
equivalence class to which $b$ projects.

\emph{Role of S2.}
S2 is the condition that ensures $\mathfrak{R}$ is $c$-uniformly
stable for all bounds arising in Lemma~\ref{lem:trunc_universal}:
the prefactor $A(K,c)$ is $c$-admissible (hence $c$-uniformly bounded),
so $b(K^*(c),c)\sim_{\mathcal{S}}\phi_{s,t}$ holds with
ratio bounds independent of all external parameters,
and $\mathfrak{R}$ maps each bound to a fixed equivalence class for all $c$.
Without S2, $A(K,c)$ may grow with $c$; the ratio
$b(K^*(c),c)/\phi_{s,t}(c)$ is not $c$-uniformly bounded;
$\mathfrak{R}(b)$ is not well-defined as a single equivalence class.
\end{definition}

\begin{remark}[Preorder vs.\ realization map: strict separation]
\label{rem:scale_admissibility}
The preorder $\succ$ lives on the \emph{target} of $\mathfrak{R}$:
it compares index pairs $(s_1,t_1)$ vs.\ $(s_2,t_2)$
and is unconditionally well-defined.
S2 governs the \emph{source} side of $\mathfrak{R}$:
it ensures that each analytic bound in $\mathcal{B}$ has a
$c$-uniformly stable image under $\mathfrak{R}$.
The two layers are logically independent:
one can have a well-defined preorder $\succ$ on $\mathcal{S}$
and simultaneously have analytic bounds whose realization map
$\mathfrak{R}$ is not $c$-uniformly stable
(precisely the situation without S2).
\end{remark}

\begin{proposition}[Scale algebra]
\label{prop:scale_algebra}
$\phi_{s_1,t_1}\cdot\phi_{s_2,t_2} = \phi_{s_1+s_2,\,t_1+t_2}$;
if $f\succ g$ then $f+g \sim_{\mathcal{S}} f$.
\end{proposition}

\begin{proposition}[Layer dominance]
\label{prop:asymp_order}
With $K=K^*(c)$, $n\geq 2$:
$\phi_{n\beta,\,1+n\gamma} \prec \phi_{\beta,\,1+\gamma}$.
\end{proposition}

%% ============================================================
\section{Logical Dependency DAG}
\label{sec:dag_logical}
%% ============================================================

\begin{remark}[DAG: organizational map and role classification]
\label{rem:dag_reading}
The diagram is an organizational dependency map; acyclicity describes
proof discipline.
The three structural axes enter the DAG at distinct levels:
\begin{itemize}[nosep]
  \item \textbf{S1 (H1):} exponent generation (lem:trunc\_univ).
  \item \textbf{S2 (EM):} ensures $c$-uniform stability of
    $\mathfrak{R}$ (Definition~\ref{def:realization_map});
    enters at rem:EM\_growth.
    Lemma~\ref{lem:S2_knockdown} shows what fails without S2:
    not $\succ$, but $\mathfrak{R}$.
  \item \textbf{S3 (subst):} diagonal projection (prop:subst).
\end{itemize}
Terminal node \textbf{thm:obstruction} receives input from
cor:rate, lem:decomp\_stab, lem:balance\_fixpt,
lem:S2\_knockdown, lem:H1\_rigidity, and three scope remarks.
\end{remark}

\[
\begin{tikzcd}[row sep=2.0em, column sep=1.1em,
  every arrow/.append style={-Latex},
  every label/.append style={font=\small}]
\boxed{\mathrm{S1/H1}} \ar[dr] \ar[drrr] \ar[drrrrrr]
& & \boxed{\mathrm{H2}} \ar[dl]
& & \boxed{F(K;c)} \ar[d]
& & & \\
& \text{lem:trunc\_univ} \ar[d] \ar[drr]
& & & \text{lem:kstar} \ar[d]
& & & \\
\boxed{\mathrm{S2/EM}} \ar[ur] \ar[rrrrrr]
& \text{thm:rate} \ar[dr]
& & \text{rem:EM\_growth} \ar[dr]
& \text{lem:loglog} \ar[d] \ar[dr]
& & \text{lem:H1\_rig} \ar[ddl]
& \text{lem:S2\_kd} \ar[ddll] \\
& & \boxed{\mathrm{S3/subst}} \ar[d] \ar[drr]
& & \text{lem:kopt(i)} \ar[d]
& \text{lem:stab} & & \\
& & \text{lem:diag\_sub} \ar[d] \ar[drr]
& & \text{lem:kopt(ii)} \ar[dl] & & & \\
& & \text{cor:rate} \ar[dr]
& & \text{lem:decomp\_stab} \ar[d] & & & \\
& & & \text{lem:balance\_fixpt} \ar[d]
& \ar[dl] & & & \\
& \text{rem:non\_taut} \ar[rr]
& & \boxed{\textbf{thm:obstruction}}
& \ar[l] \text{rem:meth\_vs\_op} & & & \\
& & \text{rem:const\_dep} \ar[urr] & & & & & \\
\end{tikzcd}
\]

\begin{remark}[Acyclicity, S2 position, obstruction node]
\label{rem:dag_nocycle}
\begin{itemize}[nosep]
  \item Anti-circularity:
    $F'\text{ sign}\to\text{window}\to\text{convexity}$, strictly directed.
  \item \textbf{S2/EM} enters in two places:
    (a) predecessor of lem:trunc\_univ
      (certifies $c$-admissibility of coefficients, ensuring
      $\mathfrak{R}$ maps each bound to a well-defined class);
    (b) direct input to thm:obstruction via lem:S2\_kd
      (certifies $c$-uniform stability of $\mathfrak{R}$ globally).
  \item lem:H1\_rigidity and lem:S2\_knockdown certify
    exponent-independence and $c$-uniform stability of $\mathfrak{R}$
    respectively; both are sibling inputs to thm:obstruction.
  \item All $K,c$-couplings pass through prop:subst (S3).
\end{itemize}
\end{remark}

%% ============================================================
\section{Framework}
\label{sec:framework}
%% ============================================================

\begin{definition}[$\mathfrak{T}_\gamma$]
\label{def:opclass}
$\mathfrak{T}_\gamma$ ($\gamma,\beta>0$): self-adjoint
$\cAg=-\partial_\xi^2+V(\xi)$ on $L^2(\R_+)$,
$V\geq 0$, $V\to+\infty$, discrete spectrum;
(i) $a_j\sim C_\gamma j^\gamma$;
(ii) $V\in C^2$, simple turning points; (iii) H1; (iv) H2.
$c$ is an approximation scale parameter, independent of the operator.
\end{definition}

\begin{hypothesis}[H1]
\label{hyp:langer_pointwise}
$\exists\,C_0<\infty$:
$\|e_j\|_{C^0([0,M])}\leq C_0\,j^\gamma c^{-\beta}$.
\end{hypothesis}

\begin{hypothesis}[H2]
\label{hyp:Cinfty}
$C_\infty=\sup_{j}\|u_j\|_{C^0([0,M])}<\infty$.
\end{hypothesis}

%% ============================================================
\section{Operator Decomposition}
\label{sec:opdecomp}
%% ============================================================

\begin{remark}[$\BK$: $c$-independence]
\label{rem:BK_definition}
$\BK$ is a functional of $\{P_j:j>K\}$ alone; no $c$ enters.
$C_2(\alpha)=\sup_{K\geq 0}\|\BK\|^2 e^{\alpha K}$
is $c$-admissible by definition.
$\|\BK\|\leq C_2^{1/2}e^{-\alpha K/2}$ holds pointwise in $K$.
\end{remark}

\begin{lemma}[$\gamma$-Agnostic nodes]
\label{lem:gamma_agnostic}
Without H1--H2:
(i) $\SK-\cKA=\AK-\BK$;
(ii) tensor norm; (iii) Bessel tail; (iv) Tonelli/DCT;
(v) $\|\BK\|\leq C_2(\alpha)^{1/2}e^{-\alpha K/2}$;
(vi) Kato.
\end{lemma}

%% ============================================================
\section{Truncation Geometry and S2 Knockdown}
\label{sec:trunc}
%% ============================================================

\begin{lemma}[Universal truncation geometry]
\label{lem:trunc_universal}
Under H1--H2, for all $K\geq 1$, $c\geq c_0$:
\begin{align}
  \normh{\AKlin}{}
    &\leq C_{\rm lin}K^{1+\gamma}c^{-\beta}
    +C_{\sigma,1}K^\gamma c^{-\beta},
    \label{eq:AKlin}\\
  \normh{\AKquad}{}
    &\leq C_{\rm quad}K^{1+2\gamma}c^{-2\beta}
    +C_{\sigma,2}K^{2\gamma}c^{-2\beta},
    \label{eq:AKquad}
\end{align}
all constants $c$-admissible.
\end{lemma}

\begin{proof}
Euler--Maclaurin:
\begin{equation}\label{eq:EM}
  \sum_{j=1}^K j^\sigma
  =\frac{K^{1+\sigma}}{1+\sigma}+\frac{1}{2}K^\sigma+R_\sigma(K),
\end{equation}
\begin{equation}\label{eq:EM_remainder}
  |R_\sigma(K)|\leq C_\sigma K^{\sigma-1},\quad
  C_\sigma=\tfrac{|B_2|}{12}\sigma(\sigma-1),\quad B_2=\tfrac{1}{6}.
\end{equation}
Multiply by $c$-admissible prefactors; close by Lemma~\ref{lem:cadm_closure}.
\end{proof}

\begin{remark}[EM structural property: two roles]
\label{rem:EM_growth}
The Euler--Maclaurin formula serves two distinct functions:
\begin{enumerate}[nosep,label=(\roman*)]
  \item \textbf{Analytic:}
    extracts the leading power $K^{1+\sigma}/(1+\sigma)$
    from the discrete spectral sum $\sum_{j=1}^K j^\sigma$,
    with remainder $O(K^{\sigma-1})$.
  \item \textbf{$c$-uniform stability of $\mathfrak{R}$:}
    the remainder coefficient $C_\sigma = \frac{|B_2|}{12}\sigma(\sigma-1)$
    is a universal explicit constant, independent of $K$ and $c$.
    This makes all prefactors in~\eqref{eq:AKlin}--\eqref{eq:AKquad}
    $c$-admissible, so each bound has a well-defined,
    $c$-uniformly stable image under $\mathfrak{R}$
    (Definition~\ref{def:realization_map}).
\end{enumerate}
Role~(ii) is structural component (S2).
EM is a canonical realization of (S2) in the WKB power-sum class;
other methods satisfying leading-power extraction,
uniform remainder control, and scaling invariance of coefficients
serve equally.
Lemma~\ref{lem:S2_knockdown} shows that without (ii),
$\mathfrak{R}$ loses $c$-uniform stability.
\end{remark}

\begin{lemma}[S2 knockdown: failure of $c$-uniform stability of $\mathfrak{R}$]
\label{lem:S2_knockdown}
Assume S2 is replaced by a summation method yielding only
\begin{equation}\label{eq:S2weak}
  \sum_{j=1}^K j^\gamma
  = a_\gamma K^{1+\gamma} + R_\gamma^*(K),
  \qquad a_\gamma\in\R_{>0},
  \qquad R_\gamma^*(K) = O_*(K^\gamma),
\end{equation}
where the prefactor $A(K,c)$ absorbing $R_\gamma^*$ is not
$c$-admissible (may grow with $c$ or depend on $K$).
Then:
\begin{itemize}[nosep]
  \item[\textup{(i)}] \textbf{$\mathfrak{R}$ is not well-defined.}
    The bound $\normh{\AKlin}{}\leq A(K,c)\,K^{1+\gamma}c^{-\beta}$
    does not have a fixed image under $\mathfrak{R}$:
    for a sequence $c_n\to\infty$ with $A(K,c_n)\to\infty$,
    the ratio $A(K,c_n)K^{1+\gamma}c_n^{-\beta}/\phi_{\beta,1+\gamma}(c_n)
    = A(K,c_n)$ is unbounded,
    so $\mathfrak{R}(b)\neq [\phi_{\beta,1+\gamma}]$ along this sequence.

  \item[\textup{(ii)}] \textbf{Comparison of bounds $\Phi_1,\Phi_2$
    via $\succ$ loses $c$-uniformity.}
    The preorder $\succ$ on $\mathcal{S}$ is defined on index pairs
    and remains unconditionally well-defined.
    However, applying $\succ$ to the \emph{bounds} $\Phi_1$ and $\Phi_2$
    requires each bound to have a $c$-uniformly stable image under
    $\mathfrak{R}$; by~(i), $\Phi_1$'s image is not stable.
    Hence the statement ``$\Phi_1\succ\Phi_2$ as bounds'' is not
    $c$-uniform: it may hold for some $c$ and fail for others
    as $A(K,c)$ grows.

  \item[\textup{(iii)}] \textbf{Balance fixed-point loses its basis.}
    Lemma~\ref{lem:balance_fixpoint}(iii) establishes
    $\Psi \asymp c^{-\beta}(\log c)^{1+\gamma}$
    via Proposition~\ref{prop:substitution_stability}(ii),
    which requires $\mathfrak{R}(\Phi_1)=[\phi_{\beta,1+\gamma}]$
    to hold $c$-uniformly.
    Without S2 this fails by~(i);
    the $\asymp$-conclusion of Lemma~\ref{lem:balance_fixpoint}(iii)
    has no valid basis.
\end{itemize}

\emph{Upshot.}
The preorder $\succ$ on $\mathcal{S}$ is not affected by the absence of S2:
it is defined on index pairs and remains unconditionally well-defined.
What S2 protects is the $c$-uniform stability of $\mathfrak{R}$:
S2 ensures that every analytic bound from
Lemma~\ref{lem:trunc_universal} has a fixed,
$c$-uniformly stable image under $\mathfrak{R}$,
and hence a well-defined position in the preorder.
The quantifier shift is precise:
with S2: $\forall K,c\;\exists C_\sigma^{\rm EM}$
($c$-independent): $|R_\sigma(K)|\leq C_\sigma^{\rm EM} K^{\sigma-1}$;
without S2: $\forall K,c\;\exists A(K,c)$
($c$-dependence not excluded): $|R^*(K)|\leq A(K,c) K^\gamma$.
Existence of $A(K,c)$ does not imply $c$-uniform stability of $\mathfrak{R}$.
\end{lemma}

\begin{proof}
\textup{(i)} By~\eqref{eq:S2weak}, the prefactor absorbing
$R_\gamma^*(K)$ is $A(K,c)$; $c$-independence is not guaranteed.
For any sequence with $A(K,c_n)\to\infty$ the ratio
$b/\phi_{\beta,1+\gamma}$ is unbounded, contradicting
$\sim_{\mathcal{S}}$ (Definition~\ref{def:scale}).

\textup{(ii)} Applying $\succ$ to bounds requires each bound
to have a fixed representative in $\mathcal{S}/\!\sim_{\mathcal{S}}$
(Definition~\ref{def:realization_map});
by~(i) $\Phi_1$'s representative shifts with $c$.

\textup{(iii)} The $\asymp$-statement requires
$c$-uniform upper and lower bounds;
the upper bound requires $\mathfrak{R}(\Phi_1)=[\phi_{\beta,1+\gamma}]$
held $c$-uniformly, which fails by~(i).
\end{proof}

%% ============================================================
\section{Variational Optimal Truncation}
\label{sec:kopt}
%% ============================================================

\begin{lemma}[$K^*(c)\to\infty$]
\label{lem:kstar_to_infty}
$K^*(c)\to\infty$ as $c\to\infty$.
\end{lemma}

\begin{lemma}[Sandwich]
\label{lem:loglog}
With $c_1=2\beta/\alpha$, $c_2=2(\beta+1)/\alpha$:
\begin{equation}\label{eq:sandwich}
  c_1\log c\leq K^*(c)\leq c_2\log c,
  \qquad \log K^*(c)=\log\log c+O(1).
\end{equation}
\end{lemma}

\begin{lemma}[Variational $K^*$]
\label{lem:kopt_variational}
\textup{(i)} On $[c_1\log c,c_2\log c]$, $F''(K;c)>0$;
unique $K^*(c)$ with
\begin{equation}\label{eq:FOC}
  L(1+\gamma){K^*}^\gamma c^{-\beta}=\tfrac{\alpha}{2}Te^{-\alpha K^*/2}.
\end{equation}

\textup{(ii)} Iterated asymptotics:
\begin{equation}\label{eq:kopt_iterated}
  K^*(c)=\frac{2\beta}{\alpha}\log c+\frac{2\gamma}{\alpha}\log\log c+O(1).
\end{equation}

\textup{(ii')} $\lambda(c)=\frac{2\gamma}{\alpha K^*(c)}\to 0$:
self-consistency certificate only.

\textup{(iii)} $F(K^*;c)=O(c^{-\beta}(\log c)^{1+\gamma})$.
\end{lemma}

\begin{lemma}[Stability]
\label{lem:stability_kopt}
$O(1)$ perturbation of~\eqref{eq:FOC} shifts $K^*$ by $O(1)$.
\begin{proof}
$\frac{dK^*}{dC_1}=\frac{1}{\alpha/2-\gamma/K^*}=O(1)$.
\end{proof}
\end{lemma}

%% ============================================================
\section{Diagonal Substitution}
\label{sec:diagonal}
%% ============================================================

\begin{lemma}[Diagonal substitution]
\label{lem:diagonal_sub}
Applies Proposition~\ref{prop:substitution_stability} at $K(c)=K^*(c)$.

$h(c):=K^*(c)/(\tfrac{2\beta}{\alpha}\log c)\to 1$;
$|h(c)-1|=O(\log\log c/\log c)$.
For large $c$: $h(c)\in[1/2,3/2]$;
constant by MVT bounded by $(1+\gamma)2^\gamma$,
finite on every compact $[\gamma_0,\gamma_1]$;
$(1+\gamma)2^\gamma\to\infty$ for $\gamma\to\infty$ explicit.
Hence
\begin{equation}\label{eq:Kstar_power}
  (K^*(c))^{1+\gamma}
  =\Bigl(\tfrac{2\beta}{\alpha}\log c\Bigr)^{1+\gamma}
  \Bigl(1+O\Bigl(\tfrac{\log\log c}{\log c}\Bigr)\Bigr).
\end{equation}
Conclusion: $\normh{\cK-\cKA}{}=O(c^{-\beta}(\log c)^{1+\gamma})$.
\end{lemma}

%% ============================================================
\section{Layer-Survival}
\label{sec:layersurvival}
%% ============================================================

\begin{theorem}[Layer-Survival]
\label{thm:layer_survival}
For $n\geq 1$, H1, all $K\geq K_0$:
$\normh{A_K^{(n)}}{}\leq C_n K^{1+n\gamma}c^{-n\beta}$.
Diagonal: $O(c^{-n\beta}(\log c)^{1+n\gamma})\prec c^{-\beta}(\log c)^{1+\gamma}$
for $n\geq 2$.
\end{theorem}

%% ============================================================
\section{Main Rate Theorem}
\label{sec:main}
%% ============================================================

\begin{theorem}[Two-scale rate]
\label{thm:universal_rate}
For each fixed $\gamma>0$, $\cAg\in\mathfrak{T}_\gamma$,
all $K\geq K_0$, $c\geq c_0$:
\[
  \normh{\cK-\cKA}{}
  \leq C_{\rm lin}K^{1+\gamma}c^{-\beta}
  +C_{\rm quad}K^{1+2\gamma}c^{-2\beta}
  +C_2(\alpha)^{1/2}e^{-\alpha K/2},
\]
all constants $c$-admissible.
\end{theorem}

\begin{corollary}[Universal logarithmic rate]
\label{cor:universal_lograte}
For each fixed $\gamma>0$, at $K=\lfloor K^*(c)\rceil$:
$\normh{\cK-\cKA}{}=O(c^{-\beta}(\log c)^{1+\gamma})$.
\end{corollary}

%% ============================================================
\section{H1 Exponent Rigidity}
\label{sec:H1rigidity}
%% ============================================================

\begin{lemma}[H1 exponent rigidity]
\label{lem:H1_rigidity}
Let $\widetilde{\mathrm{H1}}$ be any hypothesis of the form
$\|e_j\|_{C^0} \leq \tilde C_0 j^{\tilde\gamma} c^{-\tilde\beta}$,
$\tilde C_0$ $c$-admissible.

\begin{itemize}[nosep]
  \item[\textup{(i)}] \textbf{Exponent propagation.}
    The FOC~\eqref{eq:FOC} with $(\tilde\gamma,\tilde\beta)$ yields
    $\widetilde K^*(c) = \frac{2\tilde\beta}{\alpha}\log c
    + \frac{2\tilde\gamma}{\alpha}\log\log c + O(1)$
    and rate $O(c^{-\tilde\beta}(\log c)^{1+\tilde\gamma})$.
    The map $(\tilde\gamma,\tilde\beta)\mapsto(1+\tilde\gamma,\tilde\beta)$
    is injective; distinct exponent pairs yield distinct rates.

  \item[\textup{(ii)}] \textbf{Logarithmic perturbation instability.}
    If $\|e_j\|_{C^0} \leq C_0 j^\gamma (\log j)^\delta c^{-\beta}$,
    $\delta>0$, then the rate becomes
    $O(c^{-\beta}(\log c)^{1+\gamma}(\log\log c)^\delta)$,
    strictly larger than the baseline.
\end{itemize}
\end{lemma}

\begin{proof}
(i) Log-form of~\eqref{eq:FOC} is shape-invariant;
iterated asymptotics applies verbatim.
(ii) $\log K^*(c) = \log\log c + O(1)$;
pointwise substitution and Proposition~\ref{prop:scale_algebra}.
\end{proof}

\begin{remark}[Scope of H1 rigidity]
\label{rem:H1_rigidity_scope}
Status: (a)--(b) within power-law class: theorem.
(c) Envelope minimality: open (Conjecture~\ref{conj:H1_envelope}).
(d) Outside power-law class: open.
\end{remark}

%% ============================================================
\section{Structural Obstruction to Rate Improvement}
\label{sec:obstruction}
%% ============================================================

\begin{lemma}[Decomposition stability]
\label{lem:decomp_stability}
With $\Phi_1,\Phi_3\geq 0$ and $c$-admissible prefactors
(S2 ensures $\mathfrak{R}(\Phi_1)$ is well-defined;
Definition~\ref{def:realization_map}):
$\Phi_1+\Phi_3\geq\Phi_1$; no cancellation between $\Phi_1$ and $\Phi_3$.
\end{lemma}

\begin{remark}[Role]
\label{rem:decomp_role}
Rate improvement requires reducing each dominant term individually;
under (S1)--(S3) this forces the analysis in Lemma~\ref{lem:balance_fixpoint}.
\end{remark}

\begin{lemma}[Balance fixed-point]
\label{lem:balance_fixpoint}
Define $\Psi := \Phi_1+\Phi_3$.
Under \textnormal{(S1)--(S3)}:
\begin{itemize}[nosep]
  \item[\textup{(i)}] $K(c)=o(\log c)$:
    $\Psi \succ c^{-\beta}(\log c)^{1+\gamma}$.
  \item[\textup{(ii)}] $K(c)\geq C\log c$, $C>c_2$:
    $\Psi \succ c^{-\beta}(\log c)^{1+\gamma}$.
  \item[\textup{(iii)}] $K(c)\in[c_1\log c,c_2\log c]$:
    $\Psi \asymp c^{-\beta}(\log c)^{1+\gamma}$;
    no $K(c)$ in the window achieves $o(\cdot)$.
    \emph{(Requires $c$-uniform stability of $\mathfrak{R}(\Phi_1)$,
    i.e.\ S2; fails without it by Lemma~\ref{lem:S2_knockdown}(iii).)}
  \item[\textup{(iv)}] For every $K(c)$,
    $\Psi \asymp c^{-\beta}(\log c)^{1+\gamma}$ or strictly larger.
\end{itemize}
\end{lemma}

\begin{proof}
(i)--(ii): regime analysis.
(iii): $\asymp$ from~\eqref{eq:Kstar_power} via
Proposition~\ref{prop:substitution_stability}(ii);
$c$-uniform stability of $\mathfrak{R}$ certified by S2
(Remark~\ref{rem:EM_growth}(ii) and Lemma~\ref{lem:cadm_closure}).
(iv): exhaustion.
\end{proof}

\begin{remark}[Nature of the fixed point]
\label{rem:fixpoint_nature}
The balance point is an exact invariant of (S1,S2,S3), not a local minimum.
The scale $c^{-\beta}(\log c)^{1+\gamma}$ is \emph{forced}
by the three-way constraint geometry.
\end{remark}

\begin{theorem}[Structural Obstruction]
\label{thm:rate_obstruction}
Let $\cAg\in\mathfrak{T}_\gamma$ satisfy H1--H2 and~(S2),
admitting decomposition~\eqref{eq:AKlin}--\eqref{eq:AKquad}
with $c$-admissible constants.
For any truncation-based method in $\mathcal{M}(\mathrm{S1,S2,S3})$,
if $\normh{\cK-\cKA}{} = o(c^{-\beta}(\log c)^{1+\gamma})$
for some $K(c)$, then at least one of \textnormal{(S1)--(S3)} must fail.
\end{theorem}

\begin{proof}
Step~1: Lemma~\ref{lem:decomp_stability}.
Step~2: Lemma~\ref{lem:balance_fixpoint}(i)--(ii).
Step~3: Lemma~\ref{lem:balance_fixpoint}(iii)
  (uses $c$-uniform stability of $\mathfrak{R}$, guaranteed by S2).
Step~4: exhaustion.
\end{proof}

\begin{remark}[Scope]
\label{rem:obstruction_scope}
Status: conditional within $\mathcal{M}(\mathrm{S1,S2,S3})$;
not a lower bound for $\mathcal{M}^c$ (O3a: open).
\end{remark}

\begin{remark}[Non-tautology]
\label{rem:non_tautology}
$\Phi_1$ (H1-driven) and $\Phi_3$ ($\BK$-driven)
have disjoint analytic origins.
Changing H1 shifts $\Phi_1$; $\Phi_3$ is invariant.
The balance is a genuine coupling of two independent mechanisms.
\end{remark}

\begin{remark}[Constant dependence]
\label{rem:constant_dependence}
All $\asymp$-constants are $c$-uniform;
dependence on $(\gamma,\beta,\alpha,C_{\rm lin},C_2)$ is explicit.
No $K(c)$-dependence after substitution (S3 pointwise).
No $\gamma\to\infty$ uniformity claimed.
\end{remark}

\begin{remark}[Method class vs.\ operator class]
\label{rem:method_vs_operator}
$\mathcal{M}(\mathrm{S1,S2,S3})$ classifies the method, not $\mathfrak{T}_\gamma$.
(S1)--(S3) are the Tauberian condition;
$c^{-\beta}(\log c)^{1+\gamma}$ is the forced consequence.
O3a (outside $\mathcal{M}$): open.
\end{remark}

%% ============================================================
\section{Model Family}
\label{sec:models}
%% ============================================================

\begin{center}
\renewcommand{\arraystretch}{1.4}
\begin{tabular}{l l c c c c}
\toprule
\textbf{Model} & $V$ & $\gamma$ & $\beta$ & $1{+}\gamma$ & H1 status \\
\midrule
Airy     & $\xi$   & $2/3$ & $1/3$ & $5/3$ & theorem \\
Harmonic & $\xi^2$ & $1$   & $1/2$ & $2$   & conditional (sketch) \\
$|\xi|^\nu$ & & $\frac{2\nu}{\nu+2}$ & $\frac{\nu}{\nu+2}$
  & $\frac{3\nu+2}{\nu+2}$ & conditional (sketch) \\
General & & $\gamma$ & $\beta$ & $1+\gamma$ & conditional (via H1-K) \\
\bottomrule
\end{tabular}
\end{center}

%% ============================================================
\section{Open Tasks and Conjectures}
\label{sec:open}
%% ============================================================

\begin{center}
\renewcommand{\arraystretch}{1.3}
\begin{tabular}{l l l}
\toprule
\textbf{Item} & \textbf{Content} & \textbf{Status} \\
\midrule
O1 & H1 envelope minimality (Conjecture~\ref{conj:H1_envelope})
   & open (Airy: theorem) \\
O2 & H2 for $\nu\neq 1$ & open \\
O3a & lower bound for $\mathcal{M}^c$ & open \\
O3b & uniform lower bound & open (H3 not imposed) \\
O4 & Olver-Watson & open \\
\hline
\multicolumn{2}{l}{O2--O4 non-blocking for main theorems} & \checkmark \\
\bottomrule
\end{tabular}
\end{center}

\begin{conjecture}[H1 envelope minimality --- O1]
\label{conj:H1_envelope}
There exists no hypothesis
$\|e_j\|_{C^0} \leq C_0 j^\gamma f(j) c^{-\beta}$
with $f(j)=o(1)$ monotone that yields
$\Phi_1(K^*(c),c) = O(c^{-\beta}(\log c)^{1+\gamma})$
with $c$-admissible constants.

\emph{Status.}
Sufficiency: theorem (Theorem~\ref{thm:universal_rate}).
Necessity: open.
Airy case: confirmed via Abramowitz--Stegun~\S10.4.
General $V$: open.
\end{conjecture}

\begin{remark}[O1 vs.\ O3a]
\label{rem:O1_vs_O3a}
O1: axiom-minimality within $\mathcal{M}(\mathrm{S1,S2,S3})$.
O3a: method-class question outside $\mathcal{M}$.
Neither implies the other.
\end{remark}

%% ============================================================
\section{Programme Conclusion}
\label{sec:conclusion}
%% ============================================================

\textbf{XVIII--XIX--XX.}
Universal rate, structural obstruction, H1 exponent rigidity,
and S2 knockdown: the scale $c^{-\beta}(\log c)^{1+\gamma}$
is rigid within $\mathcal{M}(\mathrm{S1,S2,S3})$.
Closure inventory:
\begin{itemize}[nosep,leftmargin=1.5em]
  \item $c$-admissible chain, four sources (Remark~\ref{rem:cadmissible});
  \item \textbf{S1:} exponent generation;
  \item \textbf{S2:} $c$-uniform stability of $\mathfrak{R}$
    (Definition~\ref{def:realization_map});
    removal destroys the image of bounds under $\mathfrak{R}$,
    not the preorder $\succ$ on $\mathcal{S}$
    (Lemma~\ref{lem:S2_knockdown}, Remark~\ref{rem:scale_admissibility});
  \item \textbf{S3:} diagonal projection;
  \item Preorder $\succ$: unconditionally well-defined on index pairs;
    $c$-uniform stability of $\mathfrak{R}$ is a separate condition (S2);
  \item Decomposition stability; balance fixed-point;
    H1 exponent rigidity;
  \item Non-tautology, constant uniformity, method/operator scope;
  \item O1: Airy closed, general open.
\end{itemize}

\bigskip
\begin{quote}\itshape\large
The scale $c^{-\beta}(\log c)^{1+\gamma}$ is not merely optimal ---
it is the uniquely forced consequence of three mutually irreducible
constraints: exponent generation~(S1),
$c$-uniform stability of the realization map~(S2),
and diagonal projection~(S3).
Remove any one, and the balance fixed-point system
ceases to have a well-posed conclusion.
\end{quote}

%% ============================================================
\section*{Acknowledgements}
%% ============================================================

This paper is part of the \textit{prolate-gram-coercivity programme},
a long-term project developing spectral-geometric invariants
for WKB-class Sturm--Liouville operators.
All source files are maintained at
\texttt{https://github.com/Waschtl904/prolate-gram-coercivity}.

\begin{thebibliography}{99}

\bibitem{paper18}
S.~Schmalnauer,
\textit{Paper~XVIII: Qualitative Airy Universality and BR3 for
Spectrally Regularized WKB Operators},
manuscript, prolate-gram-coercivity programme, 2026.
\texttt{https://github.com/Waschtl904/prolate-gram-coercivity}

\bibitem{paper19}
S.~Schmalnauer,
\textit{Paper~XIX: Two-Scale Truncation Rate and Linear/Quadratic
Langer Splitting for $\mathfrak{T}_\gamma$},
manuscript, prolate-gram-coercivity programme, 2026.
\texttt{https://github.com/Waschtl904/prolate-gram-coercivity}

\bibitem{Kato1966}
T.~Kato,
\textit{Perturbation Theory for Linear Operators},
Springer, Berlin, 1966.

\bibitem{Olver1997}
F.~W.~J.~Olver,
\textit{Asymptotics and Special Functions},
AK Peters, Wellesley, 1997.

\bibitem{Zettl2005}
A.~Zettl,
\textit{Sturm--Liouville Theory},
American Mathematical Society, Providence, 2005.

\bibitem{ReedSimon2}
M.~Reed, B.~Simon,
\textit{Methods of Modern Mathematical Physics~II: Fourier Analysis,
Self-Adjointness},
Academic Press, New York, 1975.

\bibitem{Osipov2013}
A.~Osipov, V.~Rokhlin, H.~Xiao,
\textit{Prolate Spheroidal Wave Functions of Order Zero},
Springer, New York, 2013.

\bibitem{BGT1987}
N.~H.~Bingham, C.~M.~Goldie, J.~L.~Teugels,
\textit{Regular Variation},
Cambridge University Press, Cambridge, 1987.

\bibitem{AbramowitzStegun}
M.~Abramowitz, I.~A.~Stegun (eds.),
\textit{Handbook of Mathematical Functions},
National Bureau of Standards, Washington~D.C., 1964.

\end{thebibliography}

\end{document}
